% year: תשפ"ג
% semester: ב
% moed: א
% note: ד"ר פיטר סמובול
% source: exams/מבחנים/תשפג/מועד א תשפג סמסטר ב.pdf
% number: 5ב
% points: 5
% categories: antiderivatives

\begin{question}
תנו הגדרה של פונקציה קדומה לפונקציה $f(x)$. יש להוסיף דוגמה.
\end{question}

\begin{solution}
\textbf{הגדרה:} תהי $f$ מוגדרת על קטע $I$. פונקציה $F$ נקראת \textbf{פונקציה קדומה} של $f$ על $I$ אם $F$ גזירה על $I$ ומתקיים $F'(x)=f(x)$ לכל $x\in I$.
(אם $F$ קדומה של $f$ על קטע, אז כל הקדומות הן מהצורה $F+C$, $C$ קבוע.)

\textbf{דוגמה:} $F(x)=\frac{x^3}{3}$ היא פונקציה קדומה של $f(x)=x^2$ על $\R$, כי $\left(\frac{x^3}{3}\right)'=x^2$; גם $\frac{x^3}{3}+5$ קדומה שלה. דוגמה נוספת: $-\cos x$ קדומה של $\sin x$.
\end{solution}
